\section{Background and related work}
\label{sec:related}

\paragraph{Alignment and deceptive behavior.}
Learned behavior can diverge from an intended objective even when the training specification appears appropriate. Work on learned optimization distinguishes the objective used to train a system from objectives that may emerge within it \citep{hubinger2019risks}, while goal-misgeneralization experiments demonstrate that competent behavior can persist alongside pursuit of an unintended goal \citep{langosco2022goal}. These concepts motivate scrutiny of apparent alignment but do not, by themselves, establish deceptive objectives in a language model.

\paragraph{Alignment faking evaluations.}
\citet{greenblatt_alignment_faking} investigated whether a model would selectively comply with a conflicting training objective to preserve its existing preferences. \citet{koorndijk_empirical_2025} adapted a prompt-based evaluation to a small language model. Our study follows the latter's harmful-request setup and extends its scope across families, reasoning conditions, and behavioral domains. We distinguish a measurable compliance gap from the stronger claim that a model deliberately acts to preserve a future objective.

\paragraph{Evaluation signals.}
LLM judges can support large-scale response assessment \citep{zheng2023judging}, but their judgments remain dependent on the task, rubric, and evaluator. Self-reported acceptance tags are an additional signal, not a substitute for substantive response assessment. Scratchpad text can reveal references to evaluation conditions, but intermediate reasoning need not faithfully represent the computation responsible for an answer \citep{yu2025implicit}. These limitations are relevant to broader gaps in generative AI safety evaluation \citep{GapsSafetyEval2024}.

